\section{An Illustration Based on Published PHQ-ADS Results}
\label{sec:phq}

The PHQ-ADS illustrates why the fit of a merged model and the interpretation of its factor require separate assessment. Kroenke et al.\ (2016) fitted one-factor, two-factor and bifactor models to its 9 depression and 7 anxiety items in three trials of patients with chronic pain or cancer: SCOPE ($n=250$), ESCAPE ($n=241$) and INCPAD ($n=405$). The reported two-factor correlations were .912, .865 and .653, respectively. We use these correlations, the published two-factor loadings and the block lengths in a population calculation, with a common loading as a sensitivity analysis, to show how the theoretical results can inform interpretation.

The calculation assumes an exactly fitting two-factor population with the two-factor loadings reported by Kroenke et al.\ (2016, Supplemental Digital Content, Table S3). The sensitivity calculation replaces these loadings by a common loading between .70 and .80. Neither calculation is a reconstruction of the trial-specific fits. There is also an estimator difference: the published analysis treated the responses as ordered categorical in Mplus with a logistic link, whereas the population calculation applies normal-theory ML and treats the published standardized loadings as latent-response loadings. Categorical estimation can change the fit indices (Lubbe, 2023). Finally, the reported two-factor RMSEA exceeded .06 in all three trials. The population calculation therefore isolates the discrepancy due to merging the two traits under the stated assumptions; it does not allocate the observed misfit between omitted dimensionality and other sources.

\begin{table}[t]
\centering\small
\caption{Published PHQ-ADS fit indices and the population one-factor approximation. Reported values are from Kroenke et al.\ (2016), their Table 5. Population values assume an exact two-factor model with 9 depression and 7 anxiety items and the published trait correlation. Its loadings are either the published two-factor loadings (Kroenke et al., 2016, Supplemental Digital Content, Table S3) or, as a sensitivity analysis, a common loading from .70 to .80. The one-factor approximation is fitted by normal-theory ML. With the published loadings the solution is the best found by multistart numerical minimization; with a common loading it follows from the scalar classification of Appendix~\ref{app:unequal-scalar}. The last column reports the depression coefficient share in that population approximation, not the share in the published trial fit or a factor score.}
\label{tab:5}
\setlength{\tabcolsep}{4pt}
\begin{tabularx}{\textwidth}{lLLLLLLL}
\toprule
Trial & Trait correlation & Reported 1-factor RMSEA, CFI & Reported 2-factor RMSEA, CFI & Loadings & Population RMSEA & Population CFI & Depression share of the fitted factor \\
\midrule
SCOPE & .912 & .091, .956 & .085, .962 & Published & .058 & .967 & .46 \\
 & & & & Common & .037 to .060 & .969 to .981 & .58 to .59 \\
\addlinespace
ESCAPE & .865 & .083, .954 & .071, .967 & Published & .065 & .951 & .53 \\
 & & & & Common & .052 to .083 & .941 to .962 & .59 to .60 \\
\addlinespace
INCPAD & .653 & .130, .862 & .085, .941 & Published & .122 & .803 & .33 \\
 & & & & Common & .103 to .149 & .796 to .840 & .64 to .70 \\
\bottomrule
\end{tabularx}
\end{table}

The illustration shows three substantively different patterns (Table~\ref{tab:5}). At INCPAD's correlation of approximately .65, the one-factor population with the published loadings has RMSEA .122 and CFI .803, and the common-loading range gives RMSEA from .10 to .15 and CFI from .80 to .84. These values are of similar magnitude to the reported .130 and .862 and show that merging traits can itself produce substantial misfit in this range. At ESCAPE's correlation of .865, the published loadings give CFI $=.951$ but RMSEA $=.065$, the same kind of disagreement between conventional cutoffs that ESCAPE reported; a common loading near .75 gives CFI $=.953$ and RMSEA $=.065$. At SCOPE's correlation of .912, population RMSEA is .058 with the published loadings and ranges from .037 to .060 with a common loading (.0602 at loading .80, marginally above the cutoff), and CFI exceeds .96 in both calculations. The reported RMSEA is .091 for one factor and .085 for two factors, so failure of the one-factor RMSEA cutoff alone cannot attribute the discrepancy to merging depression and anxiety.

The disagreement between CFI and RMSEA follows from their different denominators. For this 16-item one-factor model, RMSEA divides the discrepancy by 104 degrees of freedom and meets .06 only when $F_{\min}\le.06^2\times104\approx.374$. CFI divides the same discrepancy by the independence-model discrepancy, which is 8.9 at ESCAPE's correlation with the published ESCAPE loadings. With these loadings, CFI reaches .95 from a trait correlation of approximately .864 and RMSEA reaches .06 from approximately .878, so ESCAPE's correlation of .865 lies just inside the interval in which the two cutoffs disagree. In the common-loading sensitivity calculation at $\varphi=.865$, the independence-model discrepancy is 7.3 for a common loading of .70, 9.4 for .75, and 12.0 for .80. With loading .70, the one-factor model meets both cutoffs from a correlation of approximately .84. With loading .75, the CFI threshold is .86 and the RMSEA threshold is .88; with loading .80, the disagreement interval extends from .88 to .912, so SCOPE's correlation lies at its upper edge. Thus, the interval widens with the loadings in this sensitivity calculation.

Which trait the fitted factor favors depends on the block lengths and the loadings together. With a common loading, the longer depression block carries more multi-information and receives a coefficient share of .58 to .59 at SCOPE's correlation and .64 to .70 at INCPAD's. The corresponding depression shares in regression or Bartlett scores are .57 to .58 and .63 to .68, compared with $9/16=.56$ in the unit-weighted sum score. With the published two-factor loadings, depression receives the larger coefficient share only in ESCAPE. In all three trials the anxiety items have higher two-factor loadings on average than the depression items, with block means of .72 and .80 in SCOPE, .71 and .75 in ESCAPE, and .66 and .75 in INCPAD, and the population one-factor approximation gives depression a coefficient share of .46, .53 and .33. When the blocks differ in length, the higher loadings alone do not determine which trait the factor favors (Section~\ref{sec:unequal}). The anxiety block has the larger total correlation in SCOPE (2.43 against 2.39) and INCPAD (2.02 against 1.70), whereas in ESCAPE the smaller loading difference leaves the larger total correlation with the longer depression block (2.13 against 1.81). With one summary loading per block, the comparison of $m_k(m_k-1)\{\lambda_k^2/(1-\lambda_k^2)\}^2$ that applies as $\varphi$ approaches 1 orders the blocks in the same way, and in these calculations the favored block is the same at every trait correlation.

The one-factor loadings implied by the published two-factor solutions agree closely with the published one-factor loadings: across the 16 items their correlations are 1.000, .999 and .961. With the published two-factor loadings as reference, the published one-factor loadings give depression coefficient shares of .46, .51 and .44. At INCPAD's correlation the regression factor score correlates .81 with depression and .94 with anxiety, whereas the unit-weighted sum correlates .89 and .87. The factor share and score share are distinct quantities, and the factor score and the sum score represent the two traits differently.

These results establish the behavior of the specified population model, not the composition of the published PHQ-ADS scores. The published loadings are estimates from a categorical model, and the reported samples add sampling uncertainty; for example, ESCAPE had $n=241$ and sample RMSEA .083 with 104 degrees of freedom. The finite-sample study provides relevant design-based context, but neither that study nor this illustration removes the estimator difference or the misspecification of the published two-factor model.
